\emph{Step 1 (two rank-one pieces).} $\Pi$ and $\J/3$ are complementary orthogonal projectors, so $M_{R}^{-1} = \Pi/\alpha + (\J/3)/M_{3}$. With $\mD = d(\J + \delta E)$, $\mD\Pi\mD^{\mathsf{T}} = d^{2}\delta^{2}E\Pi E = \tfrac23 d^{2}\delta^{2}E$, since $\J\Pi = 0$ and $e_{3}^{\mathsf{T}}\Pi e_{3} = \tfrac23$, and $\mD(\J/3)\mD^{\mathsf{T}} = \tfrac{d^{2}}{3}\,w w^{\mathsf{T}}$; this is the decomposition of Lemma \claimref{neutrino.breaking-relation} at this direction, its first term being of rank one because $E$ is. Both $\J + \delta E$ and $\J + \delta_{c}E$ annihilate $(1,-1,0)$, and both sectors commute with the exchange of the first two families, so their heavy eigenvectors (of the positive semidefinite $m_{\nu}$ and $M_{e}M_{e}^{\mathsf{T}}$) lie in the exchange-even plane, which has the orthogonal rational basis $(1,1,0)$, $e_{3}$. In that plane each heavy projector is $P = \tfrac12 + C/\sqrt{D}$ with $C$ the traceless part of the $2 \times 2$ block and $D$ its discriminant, so $|U_{\tau 3}|^{2} = \operatorname{tr}(P_{e}P_{\nu}) = \tfrac12 + \operatorname{tr}(C_{e}C_{\nu})/\sqrt{D_{e}D_{\nu}}$, in which $\operatorname{tr}(C_{e}C_{\nu})$, $D_{e}$ and $D_{\nu}$ are rational: every comparison below is an exact rational inequality after squaring with signs.

\emph{Step 2 (geometry of two rank-one pieces).} For $\lambda_{M}\hat{w}\hat{w}^{\mathsf{T}} + \lambda_{D}ee^{\mathsf{T}}$ with unit $\hat{w}, e$ at angle $\varphi$, the eigenvalues have product $\lambda_{M}\lambda_{D}\sin^{2}\varphi$ and sum $\lambda_{M} + \lambda_{D}$, so $m_{2}m_{3}/(m_{2} + m_{3})^{2} = s\sin^{2}\varphi/(1 + s)^{2}$ with $s = \lambda_{D}/\lambda_{M}$, symmetric under $s \mapsto 1/s$ and largest at $s = 1$, where the eigenvalues are $\lambda(1 \pm \cos\varphi)$ and the ratio is $\tan^{2}(\varphi/2)$. The heavy eigenvector lies in the acute sector between $\hat{w}$ and $e$, at angle $\psi$ from $\hat{w}$ with $\tan 2\psi = s\sin 2\varphi/(1 + s\cos 2\varphi)$; exchanging $s$ and $1/s$ exchanges the roles of $\hat{w}$ and $e$, so $\psi(1/s) = \varphi - \psi(s)$ and $\psi \le \varphi/2$ exactly when $s \le 1$. Since $\varphi$ depends on $\delta$ and not on $M_{3}/\alpha$, this gives (ii); $\lambda_{M} \mapsto \lambda_{D}^{2}/\lambda_{M}$ is $M_{3}/\alpha \mapsto |w|^{4}/(4\delta^{4}\,M_{3}/\alpha)$ at $d = \alpha = 1$.

\emph{Step 3 (the perturbative branch).} Write a direction in the exchange-even plane by its angle $\theta$ from $a = (1,1,0)/\sqrt{2}$ in the orthonormal basis $a, e_{3}$: $u'$ has $\tan\theta_{u} = 1/\sqrt{2}$, $\hat{w}$ has $\tan\theta_{w} = u/(3\sqrt{2})$ with $u = 3 + \delta$, and $e_{3}$ has $\theta = 90^{\circ}$; for $|\delta| \le 1$, $u \in [2, 4]$. A heavy eigenvector of a block with diagonal difference $2\Delta$ and off-diagonal entry $B'$ has $2\theta = \arg(\Delta + iB')$, and $|U_{\tau 3}|^{2} = \cos^{2}(\theta_{\nu} - \theta_{e})$. \emph{The neutrino side.} Up to the positive factor $3M_{3}/d^{2}$, $m_{\nu}$ is $s|w|^{2}E + ww^{\mathsf{T}}$ with $s = \lambda_{D}/\lambda_{M}$; its block has diagonal entries $18$ and $u^{2} + s(18 + u^{2})$ and off-diagonal entry $B' = 6u/\sqrt{2} \ge 12/\sqrt{2} > 0$, which does not depend on $s$, while $\Delta = \tfrac12(18 - u^{2} - s(18 + u^{2}))$ has $\partial_{s}\Delta = -\tfrac12(18 + u^{2}) < 0$. So $\Delta + iB'$ moves to the left along a horizontal line in the upper half-plane as $s$ grows, and $\theta_{\nu} \in (0^{\circ}, 90^{\circ})$ is strictly increasing in $s$ for every $|\delta| \le 1$ (this is also why $\psi \le \varphi/2$ exactly when $s \le 1$ in Step 2). At $s = 0$ it is $\theta_{w}$; at $s = 1$, Step 2 gives $\tan 2\psi = \sin 2\varphi/(1 + \cos 2\varphi) = \tan\varphi$ with $2\psi$ and $\varphi$ both acute, so $\psi = \varphi/2$ and $\theta_{\nu} = \tfrac12(\theta_{w} + 90^{\circ})$, the bisector of $\hat{w}$ and $e_{3}$. Hence on the branch $0 \le s \le 1$, $\theta_{w} \le \theta_{\nu} \le \tfrac12(\theta_{w} + 90^{\circ})$, and both bounds increase with $\delta$, as $\tan\theta_{w}$ does: over the whole domain $\theta_{\nu}$ lies between $\theta_{w}$ at $\delta = -1$ and the bisector at $\delta = 1$. Measured from $u'$, $\tan(\theta_{u} - \theta_{w}|_{\delta = -1}) = \sqrt{2}/8$, an angle in $(10.02^{\circ}, 10.03^{\circ})$ away from $e_{3}$, and $\tan(\theta_{w}|_{\delta = 1} - \theta_{u}) = \sqrt{2}/10$ with $90^{\circ} - \theta_{u} = \arctan\sqrt{2}$, so the neutrino tilt towards $e_{3}$ is at most $\tfrac12(\arctan(\sqrt{2}/10) + \arctan\sqrt{2}) \in (31.39^{\circ}, 31.40^{\circ})$. \emph{The charged side.} $M_{e}M_{e}^{\mathsf{T}} = c^{2}(\J + \delta_{c}E)^{2}$ has the eigenvectors of $\J + \delta_{c}E$, whose block $\bigl(\begin{smallmatrix} 2 & \sqrt{2} \\ \sqrt{2} & 1 + \delta_{c} \end{smallmatrix}\bigr)$ has both eigenvalues positive (determinant $2\delta_{c}$), so $\tan 2\theta_{e} = 2\sqrt{2}/(1 - \delta_{c})$ against $\tan 2\theta_{u} = 2\sqrt{2}$ for $\J$, and exactly $\tan 2(\theta_{e} - \theta_{u}) = 286\sqrt{2}/4357$ at $\delta_{c} = 143/500$: the charged tilt is towards $e_{3}$ and lies in $(2.65^{\circ}, 2.66^{\circ})$. \emph{The bound.} So $\theta_{\nu} - \theta_{e} \in [-12.69^{\circ}, 28.75^{\circ}]$, inside $(-90^{\circ}, 90^{\circ})$, where $\cos^{2}$ decreases with the modulus; $|U_{\tau 3}|^{2}$ is therefore least where $\theta_{\nu}$ is largest, at $\delta = 1$ and $\lambda_{M} = \lambda_{D}$, the separation there being in $(28.73^{\circ}, 28.75^{\circ})$, and Step 1's rational test puts that value in $(0.7687, 0.7688)$. This proves (iii) on the whole branch with the infimum attained. The check certifies each ingredient in exact rationals, the angles by rational enclosures of $\sqrt{2}$, of the arctangent series and of $\pi$; it confirms the block formulas exactly against $3r\,m_{\nu}$, and, independently of the monotonicity, certifies $|U_{\tau 3}|^{2} > 0.7687$ on all of $|\delta| \le 1$, $0 \le s \le 1$ by adaptive bisection, bounding $\operatorname{tr}(C_{e}C_{\nu})$ below and $D_{\nu}$ above on each box. At the instance $\delta = 0.286$, $M_{3}/\alpha = 66$ it confirms $|U_{\tau 3}|^{2} > 0.96$ exactly, together with its ratio $m_{2}/m_{3} \in (0.168, 0.172)$.

\emph{Step 4 (past the peak).} The partners of the roots $M_{3}/\alpha = 66$ ($\delta = +0.286$) and $45.3$ ($\delta = -0.286$) under Step 2's exchange are $469.5$ and $530.7$; the check confirms in exact arithmetic that each has $m_{2}/m_{3} \in (0.168, 0.172)$ and $m_{1} = 0$, that the first gives $|U_{\tau 3}|^{2} \in (0.55, 0.58)$, just outside the measured window $(0.39, 0.55)$ of the global fit \cite{NuFIT2024} (the window Observation \claimref{neutrino.mixing-none-leading-order} reads from the record; the window is data, not a premise), and the second $|U_{\tau 3}|^{2} \in (0.50, 0.55)$, inside it, and, in Step 2's terms, that partner roots have equal mass ratios and mirror-image heavy eigenvectors and that the heavy eigenvector is nearer $\hat{w}$ exactly when $\lambda_{M} > \lambda_{D}$ (400 instances), and the peak value $\tan^{2}(\varphi/2)$ with its location at $\delta = 0.286$.
